% Part: first-order-logic
% Chapter: introduction
% Section: models-theories

\documentclass[../../../include/open-logic-section]{subfiles}

\begin{document}

\olfileid[tr]{fol}{int}{mod}

\olsection{Modeller ve Kuramlar}

Birinci dereceden mantığın sözdizimini ve semantiğini tanımladıktan
sonra !!{structure}s ifadelerinin ve semantik kavramların özelliklerini
incelemeye başlayabiliriz. Ayrıca !!{derivation} sistemlerini tanımlayıp
inceleyebiliriz. Bir !!{sentence}s kümesi için şunu sorabiliriz: Bu
kümedeki bütün !!{sentence}s ifadelerini hangi !!{structure}s doğru
kılar? Bir !!{sentence}s kümesi~$\Gamma$ verildiğinde bunları sağlayan
!!a{structure}~$\Struct M$,~\emph{$\Gamma$'nın modeli} olarak
adlandırılır. $\Gamma$ ile başlayıp modellerini bulmaya çalışabiliriz:
Bunlar neye benzer? Ne kadar büyük ya da küçük olmaları gerekir? Fakat
tek bir !!{structure} ya da bir !!{structure}s topluluğuyla başlayıp
hangi !!{sentence}s ifadelerinin bunlarda doğru olduğunu da sorabiliriz.
Bu !!{structure}s ifadelerini, tam olarak bunlarda doğru olan tümceler
olmaları anlamında \emph{karakterize eden} !!{sentence}s var mıdır?
Bu tür sorular \emph{model kuramının} alanına girer. Aynı zamanda
\emph{aksiyomatik yöntemin} de temelini oluştururlar: Bir
!!{structure}s topluluğunu bir kuramın aksiyomları olan bir
!!{sentence}s kümesiyle betimlemek. Bunu mümkün kılan gözlem, tam
olarak aksiyomların birinci dereceden mantıkta mantıksal olarak
gerektirdiği !!{sentence}s ifadelerinin aksiyomların bütün modellerinde
doğru olmasıdır.

Çok basit bir örnek olarak önsıralamaları ele alalım. Bir önsıralama,
bir~$A$ kümesi üzerinde hem yansımalı hem de geçişli olan~$R$
bağıntısıdır. Üzerinde iki-li $R\subseteq A\times A$ bağıntısı bulunan
bir~$A$ kümesi, tek bir iki-li bağıntı simgesi~$\Obj P$ içeren birinci
dereceden dil için !!a{structure} vermek üzere tam olarak gereken
şeydir: $\Domain M=A$ ve $\Assign{\Obj P}M=R$ koyarız. $R$ bir
önsıralama olduğundan yansımalı ve geçişlidir; bunu söyleyen birinci
dereceden mantık !!{sentence}s kümesi~$\Gamma$ bulabiliriz:
\begin{align*}
  & \lforall[\Obj v_0][\Atom{\Obj P}{\Obj v_0,\Obj v_0}]\\
  & \lforall[\Obj v_0][\lforall[\Obj v_1][\lforall[\Obj v_2][((\Atom{\Obj P}{\Obj v_0,\Obj v_1} \land \Atom{\Obj P}{\Obj v_1,\Obj v_2}) \lif \Atom{\Obj P}{\Obj v_0,\Obj v_2})]]].
\end{align*}
Bu !!{sentence}s ifadeleri, ``her~$x$ için $Rxx$'' ($R$ yansımalıdır)
ve ``$Rxy$ ile $Ryz$ olduğunda $Rxz$ de olur'' ($R$ geçişlidir)
tümcelerinin simgeleştirmeleridir. !!a{structure}~$\Struct M$ yapısının
bu iki !!{sentence}s kümesi~$\Gamma$'nın bir modeli olması, ancak ve
ancak $R$'nin (yani $\Assign{\Obj P}M$'nin)~$A$ (yani $\Domain M$)
üzerinde bir önsıralama olmasıyla denktir. Başka bir deyişle
$\Gamma$'nın modelleri tam olarak önsıralamalardır. Yalnızca~$\Obj P$
!!{predicate} simgesini içeren birinci dereceden dilde ifade edilebilen
bütün önsıralama özellikleri (yukarıdaki yansımalılık ve geçişlilik
gibi)~$\Gamma$ içindeki iki !!{sentence}s tarafından mantıksal olarak
gerektirilir; tersi de geçerlidir. Dolayısıyla~$\Gamma$'nın modelleri
hakkında ispatladığımız her şeyi bütün önsıralamalar hakkında da
ispatlamış oluruz.

Her belirli kuram ve model sınıfı için (örneğin~$\Gamma$ ve bütün
önsıralamalar), karşılık gelen birinci dereceden dilde neyin ifade
edilip neyin edilemeyeceğine dair ilginç sorular vardır. Bütün diller
ve model sınıfları için ilginç olan bazı !!{structure}s özellikleri de
vardır; bunlar !!{domain} boyutuyla ilgilidir. Örneğin herhangi bir
$n\in\PosInt$ için !!{domain} içinde tam olarak $n$ tane !!{element}s
bulunduğu her zaman ifade edilebilir. Sonsuz sayıda !!{sentence}s
kullanarak !!{domain} ifadesinin sonsuz olduğu da ifade edilebilir.
Fakat evrenin sonlu ya da !!{nonenumerable} olduğu ifade edilemez.
Birinci dereceden dillerin bu sınırlılıkları hakkındaki sonuçlar,
tıkızlık ve Löwenheim--Skolem teoremlerinin sonuçlarıdır.

\end{document}
